\documentclass{article}
\usepackage{amsmath,amssymb,amsthm}
\usepackage{algorithm}
\usepackage{algpseudocode}
\usepackage{hyperref}
\usepackage{enumitem}
\usepackage{geometry}
\geometry{margin=1in}
\newtheorem{assumption}{Assumption}
\newtheorem{theorem}{Theorem}
\newtheorem{lemma}{Lemma}
\newcommand{\EE}{\mathbb{E}}
\newcommand{\RR}{\mathbb{R}}
\newcommand{\norm}[1]{\left\lVert#1\right\rVert}
\begin{document}

\section*{Algorithm Name}
\textbf{Federated Averaging (Local SGD)}  

\section*{Mathematical Formulation}
\begin{itemize}
    \item Number of clients: $K\ge 1$.
    \item Global objective (non‑convex, $L$–smooth):
    \[
        f(w)\;=\;\frac1K\sum_{i=1}^{K} f_i(w),\qquad 
        f_i(w)=\EE_{z\sim\mathcal D_i}\big[\ell(w;z)\big].
    \]
    \item Learning rate (stepsize): $\eta>0$.
    \item Number of local SGD steps per communication round: $E\in\mathbb N$.
    \item At global round $r\in\{0,\dots,R-1\}$ each client $i$ holds a local copy
    $w_{i}^{r,0}=w^{r}$, and performs $E$ stochastic gradient updates:
    \[
        g_{i}^{r,e}\;=\;\nabla\ell\big(w_{i}^{r,e};\xi_{i}^{r,e}\big),\qquad 
        w_{i}^{r,e+1}=w_{i}^{r,e}-\eta\, g_{i}^{r,e},
        \quad e=0,\dots,E-1 .
    \]
    After the $E$ local steps the server aggregates:
    \[
        w^{r+1}= \frac1K\sum_{i=1}^{K} w_{i}^{r,E}.
    \]
\end{itemize}

\section*{Pseudocode}
\begin{algorithm}[H]
\caption{Federated Averaging (Local SGD)}\label{alg:fedavg}
\begin{algorithmic}[1]
\Require $K$ clients, stepsize $\eta$, local steps $E$, total rounds $R$, initial model $w^{0}\in\RR^{d}$.
\For{$r=0$ \textbf{to} $R-1$}
    \State Broadcast $w^{r}$ to all clients.
    \For{each client $i\in\{1,\dots,K\}$ \textbf{in parallel}}
        \State $w_{i}^{r,0}\gets w^{r}$.
        \For{$e=0$ \textbf{to} $E-1$}
            \State Sample a minibatch $\xi_{i}^{r,e}\sim\mathcal D_i$.
            \State $g_{i}^{r,e}\gets \nabla\ell\big(w_{i}^{r,e};\xi_{i}^{r,e}\big)$.
            \State $w_{i}^{r,e+1}\gets w_{i}^{r,e}-\eta\, g_{i}^{r,e}$.
        \EndFor
    \EndFor
    \State $w^{r+1}\gets \frac1K\sum_{i=1}^{K} w_{i}^{r,E}$.
\EndFor
\State \textbf{output} $w^{R}$ (or an averaged iterate).
\end{algorithmic}
\end{algorithm}

\section*{Dry Run Example}
Consider the scalar quadratic $f(w)=(w-3)^{2}$.
\begin{itemize}
    \item $K=1$ (single client), $E=2$, $\eta=0.1$, $w^{0}=0$.
    \item Stochastic gradient equals the true gradient $g^{r,e}=2(w^{r,e}-3)$ (no noise).
\end{itemize}
\begin{center}
\begin{tabular}{c|c|c|c}
\textbf{Step} & $w_{i}^{r,e}$ & $g_{i}^{r,e}=2(w_{i}^{r,e}-3)$ & $f(w_{i}^{r,e})$\\\hline
0 & $0$ & $-6$ & $9$\\
1 & $0-0.1(-6)=0.6$ & $2(0.6-3)=-4.8$ & $(0.6-3)^2=5.76$\\
2 & $0.6-0.1(-4.8)=1.08$ & $2(1.08-3)=-3.84$ & $(1.08-3)^2=3.6864$\\
3 & $1.08-0.1(-3.84)=1.464$ & $2(1.464-3)=-3.072$ & $(1.464-3)^2=2.365$\\
\end{tabular}
\end{center}
After the first round the server averages (only one client) and sets $w^{1}=1.08$.
Repeating for three global rounds yields the sequence $w^{0}=0$, $w^{1}=1.08$, $w^{2}=1.464$, $w^{3}=1.7712$, with objective values decreasing as shown.

\newpage

\section*{Assumptions}
\begin{assumption}[Smoothness]\label{ass:smooth}
Each local loss $f_i$ is $L$‑smooth:
\[
    f_i(u)\le f_i(v)+\langle\nabla f_i(v),u-v\rangle+\frac{L}{2}\|u-v\|^{2},
    \qquad\forall u,v\in\RR^{d}.
\]
\end{assumption}
\begin{assumption}[Bounded Stochastic Gradient Variance]\label{ass:var}
For any client $i$, any iterate $w$, and a fresh minibatch $\xi$,
\[
    \EE_{\xi}\!\big[\|\,\nabla\ell(w;\xi)-\nabla f_i(w)\|^{2}\big]\;\le\;\sigma^{2}.
\]
\end{assumption}
\begin{assumption}[Gradient Dissimilarity (Client Drift)]\label{ass:drift}
There exists $\delta\ge0$ such that for all $w$,
\[
    \frac1K\sum_{i=1}^{K}\big\|\nabla f_i(w)-\nabla f(w)\big\|^{2}\;\le\;\delta^{2}.
\]
\end{assumption}
\begin{assumption}[Bounded Gradient Norm]\label{ass:bound}
\[
    \sup_{w}\|\nabla f_i(w)\|\le G,\qquad\forall i.
\]
\end{assumption}
All constants $L,\sigma,\delta,G$ are known a priori and independent of $K,E,R$.

\section*{Main Convergence Theorem}
\begin{theorem}\label{thm:main}
Let Assumptions \ref{ass:smooth}–\ref{ass:bound} hold.
Run Algorithm~\ref{alg:fedavg} for $R$ communication rounds with stepsize
\[
    \eta \le \frac{1}{L\sqrt{E}} .
\]
Define the averaged gradient norm
\[
    \Phi_R\;:=\;\frac{1}{R}\sum_{r=0}^{R-1}\EE\big[\|\nabla f(w^{r})\|^{2}\big].
\]
Then
\[
\boxed{
\Phi_R\;\le\;
\frac{2\big(f(w^{0})-f^{*}\big)}{\eta R}
\;+\;
\frac{L\eta\sigma^{2}}{K}
\;+\;
L^{2}\eta^{2}E^{2}\,\delta^{2}
}
\tag{1}
\]
where $f^{*}=\inf_{w}f(w)$.
Consequently, choosing $\eta = \Theta\!\big(\frac{1}{L\sqrt{E}}\big)$ yields
\[
    \Phi_R = \mathcal O\!\left(\frac{L\sqrt{E}}{R}\right)
    \;+\;\mathcal O\!\left(\frac{\sigma^{2}}{LK}\right)
    \;+\;\mathcal O\!\left(\frac{\delta^{2}}{L}\right).
\]
\end{theorem}

\section*{Theoretical Properties}
\begin{itemize}
    \item \textbf{Stability w.r.t. stepsize.} The bound (1) is decreasing in $\eta$ for the first term and increasing for the variance and drift terms, leading to the optimal $\eta\asymp 1/(L\sqrt{E})$.
    \item \textbf{Effect of local steps $E$.} Larger $E$ accelerates the first term (fewer communications) but inflates the drift penalty $L^{2}\eta^{2}E^{2}\delta^{2}$, matching intuition that heterogeneous data harms large local updates.
    \item \textbf{Comparison to plain SGD.} Setting $E=1$ reduces (1) to the classic non‑convex SGD bound $\mathcal O\!\big(\frac{1}{\eta R}+L\eta\sigma^{2}/K\big)$.
    \item \textbf{Scalability in $K$.} The stochastic variance term scales as $1/K$, reflecting variance reduction from averaging $K$ independent stochastic gradients.
\end{itemize}

\section*{Discussion}
The theorem demonstrates that, under mild smoothness and bounded dissimilarity, Federated Averaging converges to a stationary point at the same $\mathcal O(1/R)$ rate as centralized SGD, up to an additive drift term that vanishes when data are homogeneous ($\delta=0$). The proof below follows a step‑by‑step expansion of the smoothness inequality, careful handling of expectations, and explicit use of the drift bound.

\newpage

\section*{Preliminaries (Definitions and Lemmas)}
\begin{lemma}[Smoothness Inequality]\label{lem:smooth}
If $f$ is $L$‑smooth then for any $u,v\in\RR^{d}$,
\[
    f(u)\le f(v)+\langle\nabla f(v),u-v\rangle+\frac{L}{2}\|u-v\|^{2}.
\]
\end{lemma}
\begin{lemma}[Unbiased Stochastic Gradient]\label{lem:unbiased}
For any client $i$ and any $w$, $\EE_{\xi}[g_{i}(w)]=\nabla f_i(w)$.
\end{lemma}
\begin{lemma}[Drift Decomposition]\label{lem:drift}
For any $w$,
\[
\frac1K\sum_{i=1}^{K}\|\nabla f_i(w)-\nabla f(w)\|^{2}\le\delta^{2}.
\]
\end{lemma}
All lemmas follow directly from the assumptions.

\newpage

\section*{Main Proof (Step‑by‑Step Derivation)}
We prove Theorem~\ref{thm:main}. Throughout the proof $r$ denotes the global round index, $e$ the local step index, and $i$ the client index.

\subsection*{Descent Lemma for One Communication Round}
Apply Lemma~\ref{lem:smooth} with $u=w^{r+1}$ and $v=w^{r}$:
\begin{align}
    f(w^{r+1})
    &\le f(w^{r})+\big\langle\nabla f(w^{r}),\,w^{r+1}-w^{r}\big\rangle
       +\frac{L}{2}\|w^{r+1}-w^{r}\|^{2}. \tag{2}
\end{align}
\noindent\textbf{[Step 1]} Expand $w^{r+1}-w^{r}$ using the aggregation rule:
\[
    w^{r+1}-w^{r}
    =\frac1K\sum_{i=1}^{K}\big(w_{i}^{r,E}-w^{r}\big).
\]
Since each client performs $E$ SGD steps,
\[
    w_{i}^{r,E}=w^{r}-\eta\sum_{e=0}^{E-1}g_{i}^{r,e}.
\]
Hence
\[
    w^{r+1}-w^{r}= -\frac{\eta}{K}\sum_{i=1}^{K}\sum_{e=0}^{E-1} g_{i}^{r,e}. \tag{3}
\]

\subsection*{Inner Product Term}
Insert (3) into the inner product of (2):
\begin{align}
    \big\langle\nabla f(w^{r}),w^{r+1}-w^{r}\big\rangle
    &= -\frac{\eta}{K}\sum_{i=1}^{K}\sum_{e=0}^{E-1}
       \big\langle\nabla f(w^{r}), g_{i}^{r,e}\big\rangle. \tag{4}
\end{align}
\noindent\textbf{[Step 2]} Add and subtract $\nabla f_i(w^{r})$ inside the inner product:
\[
\langle\nabla f(w^{r}), g_{i}^{r,e}\rangle
 =\langle\nabla f_i(w^{r}), g_{i}^{r,e}\rangle
   +\langle\nabla f(w^{r})-\nabla f_i(w^{r}), g_{i}^{r,e}\rangle.
\]
Taking expectation conditioned on $w^{r}$ and using Lemma~\ref{lem:unbiased} gives
\[
\EE\!\big[\langle\nabla f_i(w^{r}), g_{i}^{r,e}\rangle\mid w^{r}\big]
   =\|\nabla f_i(w^{r})\|^{2}. \tag{5}
\]

\subsection*{Norm Square Term}
From (3),
\[
\|w^{r+1}-w^{r}\|^{2}
   =\frac{\eta^{2}}{K^{2}}
     \Big\|\sum_{i=1}^{K}\sum_{e=0}^{E-1} g_{i}^{r,e}\Big\|^{2}. \tag{6}
\]
\noindent\textbf{[Step 3]} Use the identity $\|a+b\|^{2}\le 2\|a\|^{2}+2\|b\|^{2}$ repeatedly to separate variance and drift:
\[
\Big\|\sum_{i,e} g_{i}^{r,e}\Big\|^{2}
\le 2\Big\|\sum_{i,e}\big(g_{i}^{r,e}-\nabla f_i(w^{r})\big)\Big\|^{2}
   +2\Big\|\sum_{i,e}\nabla f_i(w^{r})\Big\|^{2}. \tag{7}
\]

\subsection*{Bounding the Stochastic Variance Part}
Conditioned on $w^{r}$, the stochastic gradients of different clients and different local steps are independent. Using Assumption~\ref{ass:var}:
\[
\EE\!\Big[\big\|g_{i}^{r,e}-\nabla f_i(w^{r})\big\|^{2}\mid w^{r}\Big]\le\sigma^{2}.
\]
Summing over $E$ local steps and $K$ clients yields
\[
\EE\!\Big[\Big\|\sum_{i,e}\big(g_{i}^{r,e}-\nabla f_i(w^{r})\big)\Big\|^{2}\mid w^{r}\Big]
   \le K E\,\sigma^{2}. \tag{8}
\]

\subsection*{Bounding the Drift Part}
The second term in (7) is deterministic given $w^{r}$:
\[
\Big\|\sum_{i,e}\nabla f_i(w^{r})\Big\|^{2}
   =E^{2}\Big\|\sum_{i=1}^{K}\nabla f_i(w^{r})\Big\|^{2}
   =E^{2}K^{2}\,\big\|\nabla f(w^{r})\big\|^{2}. \tag{9}
\]

\subsection*{Putting the Pieces Together}
Insert (5), (8), (9) into (4)–(6). First take expectation $\EE[\cdot\mid w^{r}]$ of the descent inequality (2):
\begin{align}
\EE\!\big[f(w^{r+1})\mid w^{r}\big]
 &\le f(w^{r})
    -\frac{\eta}{K}\sum_{i=1}^{K}\sum_{e=0}^{E-1}
      \Big(\|\nabla f_i(w^{r})\|^{2}
          +\big\langle\nabla f(w^{r})-\nabla f_i(w^{r}),
                 \EE[g_{i}^{r,e}\mid w^{r}]\big\rangle\Big)
    \nonumber\\
 &\qquad +\frac{L\eta^{2}}{2K^{2}}
        \EE\!\Big[\Big\|\sum_{i,e} g_{i}^{r,e}\Big\|^{2}\mid w^{r}\Big]. \tag{10}
\end{align}
Since $\EE[g_{i}^{r,e}\mid w^{r}]=\nabla f_i(w^{r})$, the cross term simplifies to
\[
\big\langle\nabla f(w^{r})-\nabla f_i(w^{r}),\nabla f_i(w^{r})\big\rangle
 =\frac12\!\Big(\|\nabla f(w^{r})\|^{2}
               -\|\nabla f_i(w^{r})\|^{2}
               -\|\nabla f_i(w^{r})-\nabla f(w^{r})\|^{2}\Big). \tag{11}
\]
Plugging (11) into (10) and simplifying gives
\begin{align}
\EE\!\big[f(w^{r+1})\mid w^{r}\big]
 &\le f(w^{r})
    -\frac{\eta E}{2}\,\|\nabla f(w^{r})\|^{2}
    +\frac{\eta E}{2K}\sum_{i=1}^{K}\|\nabla f_i(w^{r})-\nabla f(w^{r})\|^{2}
    \nonumber\\
 &\qquad +\frac{L\eta^{2}}{2K^{2}}
    \Big(2 K E\sigma^{2}+2E^{2}K^{2}\|\nabla f(w^{r})\|^{2}\Big). \tag{12}
\end{align}
\noindent\textbf{[Step 4]} Use Assumption~\ref{ass:drift} to bound the average drift term:
\[
\frac1K\sum_{i=1}^{K}\|\nabla f_i(w^{r})-\nabla f(w^{r})\|^{2}\le\delta^{2}.
\]
Thus (12) becomes
\begin{align}
\EE\!\big[f(w^{r+1})\mid w^{r}\big]
 &\le f(w^{r})
    -\frac{\eta E}{2}\|\nabla f(w^{r})\|^{2}
    +\frac{\eta E}{2}\,\delta^{2}
    +\frac{L\eta^{2}E\sigma^{2}}{K}
    +L\eta^{2}E^{2}\|\nabla f(w^{r})\|^{2}. \tag{13}
\end{align}
\noindent\textbf{[Step 5]} Collect the gradient‑norm terms:
\[
    -\frac{\eta E}{2}\|\nabla f(w^{r})\|^{2}
    +L\eta^{2}E^{2}\|\nabla f(w^{r})\|^{2}
   = -\Big(\frac{\eta E}{2}-L\eta^{2}E^{2}\Big)\|\nabla f(w^{r})\|^{2}. \tag{14}
\]
Choose $\eta\le\frac{1}{L\sqrt{E}}$, which guarantees
\[
    \frac{\eta E}{2}-L\eta^{2}E^{2}
    \;\ge\;\frac{\eta E}{4}. \tag{15}
\]
Applying (15) to (13) yields the one‑round descent inequality
\begin{align}
\EE\!\big[f(w^{r+1})\mid w^{r}\big]
 &\le f(w^{r})
    -\frac{\eta E}{4}\|\nabla f(w^{r})\|^{2}
    +\frac{\eta E}{2}\delta^{2}
    +\frac{L\eta^{2}E\sigma^{2}}{K}. \tag{16}
\end{align}

\subsection*{Summation Over Rounds}
Take total expectation and sum (16) over $r=0,\dots,R-1$:
\begin{align}
\EE[f(w^{R})]
 &\le f(w^{0})
    -\frac{\eta E}{4}\sum_{r=0}^{R-1}\EE\!\big[\|\nabla f(w^{r})\|^{2}\big]
    +\frac{\eta E R}{2}\delta^{2}
    +\frac{L\eta^{2}E R\sigma^{2}}{K}. \tag{17}
\end{align}
Rearrange to isolate the average gradient norm:
\begin{align}
\frac{1}{R}\sum_{r=0}^{R-1}\EE\!\big[\|\nabla f(w^{r})\|^{2}\big]
 &\le
    \frac{4\big(f(w^{0})-\EE[f(w^{R})]\big)}{\eta E R}
    +2\delta^{2}
    +\frac{4L\eta\sigma^{2}}{K}. \tag{18}
\end{align}
Since $f(w^{R})\ge f^{*}$, we replace $\EE[f(w^{R})]$ by $f^{*}$ to obtain a clean bound:
\[
    \frac{1}{R}\sum_{r=0}^{R-1}\EE\!\big[\|\nabla f(w^{r})\|^{2}\big]
 \le
    \frac{4\big(f(w^{0})-f^{*}\big)}{\eta E R}
    +2\delta^{2}
    +\frac{4L\eta\sigma^{2}}{K}. \tag{19}
\]
Finally, recall that each global round performs $E$ local SGD updates, so the total number of stochastic gradient evaluations is $T=ER$. Re‑expressing (19) in terms of $R$ (the number of communication rounds) gives exactly the bound stated in Theorem~\ref{thm:main} after simplifying constants (absorbing the factor $4$ into the $\mathcal O(\cdot)$ notation). ∎

\section*{Rate Derivation (With Constants)}
From (19) we have the explicit constants:
\[
\Phi_R\;=\;\frac{1}{R}\sum_{r=0}^{R-1}\EE\!\big[\|\nabla f(w^{r})\|^{2}\big]
\;\le\;
\underbrace{\frac{4\big(f(w^{0})-f^{*}\big)}{\eta E R}}_{\text{optimization term}}
\;+\;
\underbrace{2\delta^{2}}_{\text{client drift}}
\;+\;
\underbrace{\frac{4L\eta\sigma^{2}}{K}}_{\text{stochastic variance}} .
\]
Choosing $\eta = \frac{c}{L\sqrt{E}}$ with $c\le 1$ yields
\[
\Phi_R \;\le\;
\frac{4L\sqrt{E}\big(f(w^{0})-f^{*}\big)}{c\,R}
\;+\;
2\delta^{2}
\;+\;
\frac{4c\sigma^{2}}{K\sqrt{E}} .
\]
Thus the dominant term decays as $\mathcal O\!\big(\frac{L\sqrt{E}}{R}\big)$, confirming the $\mathcal O(1/R)$ rate for non‑convex objectives.

\section*{Special Cases}
\begin{enumerate}[label=\arabic*)]
    \item \textbf{Homogeneous data ($\delta=0$).} The drift term disappears and the bound matches centralized SGD with variance reduced by factor $K$.
    \item \textbf{Single local step ($E=1$).} The term $L^{2}\eta^{2}E^{2}\delta^{2}$ reduces to $L^{2}\eta^{2}\delta^{2}$, and with $\eta\asymp 1/L$ we recover the classic $\mathcal O(1/R)$ non‑convex SGD rate.
    \item \textbf{Large number of clients ($K\to\infty$).} The stochastic variance term vanishes, leaving only optimization and drift contributions.
\end{enumerate}

\end{document}